\section{Week 7: Evaluating Models II: Information Criteria, Parameter
Recovery and Model Recovery}

\begin{quote}
\textbf{Chapter 3: Evaluation of cognitive models} · Lecture 1 h · Lab 2
h \textbf{Running data:} the Week 2--6 forgetting data and models, plus
large simulation studies at the same design.
\end{quote}

\begin{center}\rule{0.5\linewidth}{0.5pt}\end{center}

\subsection{Key terms}

Important terms of this lesson, with their Persian equivalents (repeated
even if they appeared in earlier weeks).

{\def\LTcaptype{none} % do not increment counter
\begin{longtable}[]{@{}
  >{\raggedright\arraybackslash}p{(\linewidth - 4\tabcolsep) * \real{0.3333}}
  >{\raggedright\arraybackslash}p{(\linewidth - 4\tabcolsep) * \real{0.3333}}
  >{\raggedright\arraybackslash}p{(\linewidth - 4\tabcolsep) * \real{0.3333}}@{}}
\toprule\noalign{}
\begin{minipage}[b]{\linewidth}\raggedright
English
\end{minipage} & \begin{minipage}[b]{\linewidth}\raggedright
\RL{\CMfa فارسی}
\end{minipage} & \begin{minipage}[b]{\linewidth}\raggedright
\RL{\CMfa توضیح کوتاه}
\end{minipage} \\
\midrule\noalign{}
\endhead
\bottomrule\noalign{}
\endlastfoot
criterion (c) & \RL{\CMfa ملاک} & \RL{\CMfa آستانه‌ی تصمیم؛ نشان‌دهنده‌ی سوگیری
پاسخ} \\
parameter recovery & \RL{\CMfa بازیابی پارامتر} & \RL{\CMfa آیا برازش، پارامترهای
معلوم شبیه‌سازی را برمی‌گرداند} \\
parameter & \RL{\CMfa پارامتر} & \RL{\CMfa کمیت قابل تنظیم مدل که بین افراد یا شرایط
فرق می‌کند} \\
model recovery & \RL{\CMfa بازیابی مدل} & \RL{\CMfa آیا مقایسه‌ی مدل، مدل مولد را
درست انتخاب می‌کند} \\
forgetting curve & \RL{\CMfa منحنی فراموشی} & \RL{\CMfa کاهش یادآوری با گذشت
زمان} \\
simulation & \RL{\CMfa شبیه‌سازی} & \RL{\CMfa تولید داده‌ی مصنوعی با اجرای مدل} \\
research design & \RL{\CMfa طرح آزمایش} & \RL{\CMfa شرایط، تعداد کوشش‌ها و
نمونه‌ها} \\
likelihood-ratio test (LRT) & \RL{\CMfa آزمون نسبت درست‌نمایی} & \RL{\CMfa مقایسه‌ی دو
مدل تودرتو} \\
likelihood & \RL{\CMfa درست‌نمایی} & \RL{\CMfa احتمال داده به‌عنوان تابعی از
پارامترها} \\
nested models & \RL{\CMfa مدل‌های تودرتو} (nested) & \RL{\CMfa یک مدل حالت خاص مدل
دیگر است} \\
AIC / BIC & \RL{\CMfa معیار اطلاعاتی آکائیکه} / \RL{\CMfa بیزی} & \RL{\CMfa برازش به‌علاوه‌ی
جریمه برای تعداد پارامتر} \\
response bias & \RL{\CMfa سوگیری پاسخ} & \RL{\CMfa تمایل به یک پاسخ، مستقل از
محرک} \\
confusion matrix & \RL{\CMfa ماتریس درهم‌ریختگی} (confusion matrix) & P(\RL{\CMfa مدل
انتخاب‌شده} \\
inversion matrix & \RL{\CMfa ماتریس معکوس} (inversion matrix) & P(\RL{\CMfa مدل
مولد} \\
fit / fitting & \RL{\CMfa برازش} & \RL{\CMfa تنظیم پارامترها تا مدل با داده جور
شود} \\
exponential / power function & \RL{\CMfa تابع نمایی} / \RL{\CMfa توانی} & \RL{\CMfa دو شکل
رایج برای منحنی فراموشی و یادگیری} \\
identifiability & \RL{\CMfa شناسایی‌پذیری} & \RL{\CMfa آیا داده پارامترها را به‌طور
یکتا تعیین می‌کند} \\
parameter trade-off & \RL{\CMfa بده‌بستان پارامترها} & \RL{\CMfa تغییر یک پارامتر که
با پارامتر دیگر جبران می‌شود} \\
log-likelihood / NLL & \RL{\CMfa لگاریتم درست‌نمایی} / \RL{\CMfa منفی آن} &
\RL{\CMfa مقداری که بهینه‌سازها کمینه یا بیشینه می‌کنند} \\
Akaike weights & \RL{\CMfa وزن‌های آکائیکه} & \RL{\CMfa پشتیبانی نسبی از هر مدل} \\
cross-validation (CV) & \RL{\CMfa اعتبارسنجی متقابل} (cross-validation) &
\RL{\CMfa برازش روی بخشی از داده و آزمون روی بخش کنارگذاشته} \\
trial & \RL{\CMfa کوشش} (trial) & \RL{\CMfa یک بار انجام تکلیف توسط شرکت‌کننده} \\
model mimicry & \RL{\CMfa تقلید مدل} & \RL{\CMfa تولید داده‌ای که شبیه داده‌ی مدل دیگر
است} \\
in-context learning & \RL{\CMfa یادگیری درون‌متنی} & \RL{\CMfa یادگیری از مثال‌های
داخل} prompt \\
\end{longtable}
}

\begin{center}\rule{0.5\linewidth}{0.5pt}\end{center}

\subsection{Learning outcomes}

By the end of this week, students will be able to:

\begin{enumerate}
\def\labelenumi{\arabic{enumi}.}
\tightlist
\item
  \textbf{Conduct} a likelihood-ratio test for nested models and
  \textbf{explain} why it cannot compare non-nested models.
\item
  \textbf{Compute} and \textbf{interpret} AIC, AICc, BIC and Akaike/BIC
  weights, and \textbf{explain} the different goals of AIC and BIC.
\item
  \textbf{Design} and \textbf{run} a parameter-recovery study, and
  \textbf{diagnose} biased, noisy or traded-off parameters.
\item
  \textbf{Design} and \textbf{run} a model-recovery study, and
  \textbf{interpret} the confusion and inversion matrices.
\item
  \textbf{Evaluate} whether a given experimental design and comparison
  method can support the conclusions a researcher wants to draw.
\end{enumerate}

\begin{center}\rule{0.5\linewidth}{0.5pt}\end{center}

\subsection{Motivation \& big picture}

Cross-validation (Week 6) is general but expensive: every fold means
refitting, and refitting complex models with multi-start can take hours.
\textbf{Information criteria} estimate out-of-sample performance from a
\emph{single} fit, by adding a complexity penalty to the training fit.

But there is a more fundamental question that no criterion answers on
its own: \textbf{can our method even work at this design?} Suppose we
conclude ``participants forget by a power law, with decay exponent
0.42''. Two things must be true for this to mean anything:

\begin{enumerate}
\def\labelenumi{\arabic{enumi}.}
\tightlist
\item
  If the data had come from a power model with \(d = 0.42\), our fitting
  procedure would have returned something close to 0.42
  (\textbf{parameter recovery}).
\item
  If the data had come from an exponential model, our comparison method
  would have said ``exponential'', not ``power'' (\textbf{model
  recovery}).
\end{enumerate}

Both can be checked by simulation, because we know the truth for
simulated data. Modern guidelines for computational modeling treat these
checks as essential; skipping them is like running an experiment without
checking that the equipment works. This week completes the ``simulate \ensuremath{\rightarrow}
fit \ensuremath{\rightarrow} evaluate \ensuremath{\rightarrow} \textbf{recover}'' part of our pipeline.

\paragraph{Where is this used?}

\begin{itemize}
\tightlist
\item
  \textbf{Model comparison tables with AIC/BIC} appear in almost every
  modeling paper in psychology, neuroscience, ecology and economics.
\item
  \textbf{Computational psychiatry}: before claiming that patients have
  a lower learning rate, researchers must show by parameter recovery
  that the learning rate can be estimated at their design. Reviewers
  increasingly require this.
\item
  \textbf{Designing experiments}: recovery simulations are run
  \emph{before} collecting data, to choose the number of trials and
  conditions (a modeling version of power analysis).
\item
  \textbf{Any study that says ``model X explains the data best''} should
  report model recovery; otherwise the winner may only reflect a bias of
  the criterion.
\end{itemize}

\begin{center}\rule{0.5\linewidth}{0.5pt}\end{center}

\subsection{Lecture notes}

\paragraph{In plain words}

Cross-validation is honest but slow. \textbf{Information criteria} (AIC,
BIC) are quick formulas that take the fit of each model and subtract a
penalty for every parameter, so that a model only wins if its extra
parameters improve the fit by more than they ``cost''. AIC aims at the
best predictor, BIC at the true model, and BIC's penalty is larger. The
second, more important half of the week asks whether our methods work at
all. We test them on fake data where we \emph{know} the truth.
\textbf{Parameter recovery}: simulate data with known parameter values,
fit, and check that the fitted values come back close to the true ones.
\textbf{Model recovery}: simulate data from each candidate model and
check that the comparison method picks the model that actually generated
them. If a method cannot find the truth in simulated data, its answer on
real data means nothing.

\textbf{Suggested timing (60 min):} likelihood-ratio test 8 · AIC and
BIC 20 · parameter recovery 14 · model recovery 14 · identifiability 4.

\subsubsection{1. The likelihood-ratio test (nested models)}

Model \(\mathcal{M}_0\) is \textbf{nested} in \(\mathcal{M}_1\) if
\(\mathcal{M}_0\) is a special case of \(\mathcal{M}_1\) obtained by
fixing some parameters. Example: the exponential model is the
exponential-with-asymptote model with \(s = 0\) (asymptote at 0). Under
the null hypothesis that \(\mathcal{M}_0\) is true,

\[
\Lambda = 2\big[\ell_1(\hat{\boldsymbol\theta}_1) - \ell_0(\hat{\boldsymbol\theta}_0)\big] \;\dot\sim\; \chi^2_{\Delta k},
\]

where \(\Delta k\) is the difference in the number of free parameters. A
large \(\Lambda\) means the extra parameters improve the fit more than
chance would.

\textbf{Limitations:} it only works for \textbf{nested} models (power
vs.~exponential cannot be tested this way); it is a null-hypothesis test
(it can reject the simpler model but never support it); and it is
approximate when the fixed parameter value lies on the boundary of the
parameter space (as \(s = 0\) does).

\begin{quote}
\textbf{Advanced: boundary nulls.} When the null value is on the
boundary (a variance of 0, an asymptote of 0, a mixture weight of 0),
the null distribution of \(\Lambda\) is a 50:50 mixture of \(\chi^2_0\)
(a point mass at 0) and \(\chi^2_1\). Using \(\chi^2_1\) is then
conservative: the correct \(p\)-value is half of the naive one.
\end{quote}

\subsubsection{2. Information criteria}

Information criteria trade off fit against complexity. With maximized
log-likelihood \(\hat\ell\), \(k\) free parameters and \(N\)
observations:

\[
\mathrm{AIC} = -2\hat\ell + 2k, \qquad
\mathrm{AICc} = \mathrm{AIC} + \frac{2k(k+1)}{N - k - 1}, \qquad
\mathrm{BIC} = -2\hat\ell + k \ln N.
\]

\textbf{Lower is better.} Only \textbf{differences} between models
fitted to the same data are meaningful.

{\def\LTcaptype{none} % do not increment counter
\begin{longtable}[]{@{}
  >{\raggedright\arraybackslash}p{(\linewidth - 4\tabcolsep) * \real{0.3333}}
  >{\raggedright\arraybackslash}p{(\linewidth - 4\tabcolsep) * \real{0.3333}}
  >{\raggedright\arraybackslash}p{(\linewidth - 4\tabcolsep) * \real{0.3333}}@{}}
\toprule\noalign{}
\begin{minipage}[b]{\linewidth}\raggedright
\end{minipage} & \begin{minipage}[b]{\linewidth}\raggedright
AIC (and AICc)
\end{minipage} & \begin{minipage}[b]{\linewidth}\raggedright
BIC
\end{minipage} \\
\midrule\noalign{}
\endhead
\bottomrule\noalign{}
\endlastfoot
Goal & Best \textbf{prediction} of new data (minimize expected
Kullback--Leibler divergence to the truth) & Pick the \textbf{true}
model if it is among the candidates (consistency); approximates the
Bayes factor \\
Penalty per parameter & 2 & \(\ln N\) (bigger than 2 once
\(N \ge 8\)) \\
As \(N \to \infty\) & May keep choosing over-complex models & Chooses
the true model with probability \ensuremath{\rightarrow} 1 \\
Small-sample correction & AICc (use when \(N/k < 40\)) & --- \\
\end{longtable}
}

\textbf{What is \(N\)?} The number of \emph{independent observations}
that enter the likelihood. For binomial recall data that is the number
of item responses (here \(8 \times 50 = 400\)), not the number of cells
(8). This choice affects BIC (through \(\ln N\)) and AICc. Report it.

\textbf{Weights.} With
\(\Delta_i = \mathrm{IC}_i - \min_j \mathrm{IC}_j\),

\[
w_i = \frac{\exp(-\Delta_i/2)}{\sum_j \exp(-\Delta_j/2)}.
\]

\textbf{Akaike weights} estimate the probability that model \(i\) is the
best predictor among the candidates. \textbf{BIC weights} approximate
posterior model probabilities under equal prior odds (Week 10).

\textbf{Rules of thumb for \(\Delta\):} 0--2, essentially equivalent
support; 4--7, considerably less support; more than 10, essentially
none. These are conventions, not laws.

\textbf{Limitation shared by AIC and BIC:} both count parameters, \(k\),
so they are blind to differences in \textbf{functional-form flexibility}
between models with the same \(k\) (Week 6). Cross-validation and Bayes
factors (Week 10) are not.

\begin{quote}
\textbf{Advanced: where the AIC penalty comes from.} The maximized
log-likelihood is an optimistically biased estimate of the expected
log-likelihood of new data, because the same data were used to choose
\(\hat{\boldsymbol\theta}\). Akaike showed that, for regular models and
large \(N\), the bias is approximately \(k\). Correcting for it gives
\(-2(\hat\ell - k) = \mathrm{AIC}\). That is also why AIC and LOO-CV
agree asymptotically.
\end{quote}

\paragraph{Reading the math}

\begin{itemize}
\tightlist
\item
  \(-2\hat\ell\) is the \textbf{misfit} part (twice the minimized NLL);
  \(2k\) or \(k\ln N\) is the \textbf{penalty} part. Lower AIC/BIC is
  better.
\item
  Each extra parameter must lower \(-2\hat\ell\) by more than 2 (AIC) or
  more than \(\ln N\) (BIC) to be worth it. With \(N = 400\),
  \(\ln N \approx 6\): a parameter must improve the NLL by about 3 units
  to pass BIC, but only by 1 unit to pass AIC.
\item
  \textbf{Weights}: \(\exp(-\Delta_i/2)\) converts a difference in AIC
  into a relative likelihood. \(\Delta = 2\) gives \(e^{-1} = 0.37\), so
  a model 2 AIC units worse is about a third as well supported. Dividing
  by the sum makes the weights add up to 1.
\item
  \textbf{Worked numbers:} power model AIC
  \(= 2(19.830) + 2(2) = 43.66\); exponential-with-asymptote
  \(= 2(18.919) + 2(3) = 43.84\). The asymptote model fits better (lower
  NLL) but pays for one more parameter, so the two end up almost tied.
\end{itemize}

\subsubsection{3. Parameter recovery}

\textbf{Question:} at this design, does the fitting procedure return the
parameters that generated the data?

\textbf{Procedure:} 1. Choose a set of true parameter values
\(\boldsymbol\theta^\ast\) covering the plausible range (a grid, or
random draws from plausible ranges, e.g., around the estimates from real
data). 2. For each \(\boldsymbol\theta^\ast\), simulate a dataset with
\textbf{exactly the real design} (same conditions, numbers of trials).
3. Fit the model with \textbf{exactly the real fitting procedure} (same
optimizer, starts, bounds). 4. Compare \(\hat{\boldsymbol\theta}\) with
\(\boldsymbol\theta^\ast\): scatter plots of true vs.~recovered;
correlation; bias (mean of \(\hat\theta - \theta^\ast\)); RMSE. 5. Also
plot the recovered parameters \textbf{against each other}. A strong
correlation between, say, \(\hat\lambda\) and \(\hat s\) that is absent
in the true values reveals a \textbf{trade-off}.

\textbf{Reading the results:} points on the identity line mean good
recovery; a cloud parallel to the x-axis means the parameter is not
recoverable at this design; a curve bending away at the extremes means
bias near the bounds. If recovery fails, do not interpret that
parameter. Change the design (more trials, other conditions) or simplify
the model.

\subsubsection{4. Model recovery}

\textbf{Question:} if model \(\mathcal{M}_a\) generated the data, does
our comparison method select \(\mathcal{M}_a\)?

\textbf{Procedure:} 1. For each candidate model \(\mathcal{M}_a\):
simulate many datasets at the real design, with parameters drawn from
plausible ranges. 2. Fit \textbf{every} candidate model to every
dataset. 3. Select a winner by the chosen criterion (AIC, BIC, CV,
\ldots). 4. Tabulate the \textbf{confusion matrix}
\(P(\text{selected} = b \mid \text{generated} = a)\).

A good method gives a matrix with large diagonal entries. Off-diagonal
mass shows \textbf{model mimicry}: one model can produce data that look
like another's.

What we actually need when interpreting real data is the reverse: given
that model \(b\) won, how likely is it that \(b\) generated the data?
This is the \textbf{inversion matrix}
\(P(\text{generated} = a \mid \text{selected} = b)\), obtained by
normalizing each \emph{column} of the confusion counts (assuming equal
prior frequency of the generating models).

\begin{quote}
\textbf{Note:} model recovery also reveals \textbf{biases of criteria}:
BIC tends to favor simpler models at small \(N\), AIC more complex ones.
The ``best'' criterion is the one that recovers your candidate models
best \emph{at your design}, which is exactly what this study measures.
\end{quote}

\paragraph{Reading the math}

\begin{itemize}
\tightlist
\item
  \textbf{Confusion matrix}
  \(P(\text{selected} = b \mid \text{generated} = a)\): rows are the
  true (generating) models, columns the selected ones. Each row adds up
  to 1. Read a row as: ``when the truth is \(a\), how often do we pick
  each model?''
\item
  \textbf{Inversion matrix}
  \(P(\text{generated} = a \mid \text{selected} = b)\): the same counts,
  but each \emph{column} is normalized to 1. Read a column as: ``when we
  picked \(b\), how often was \(b\) really the truth?'' This is the
  question that matters for a real dataset, where we only see the
  winner.
\item
  \textbf{Small example.} 100 datasets from each of two models. From A:
  90 select A, 10 select B. From B: 40 select A, 60 select B. Confusion
  row A = (0.90, 0.10). But among the 130 datasets that selected A, only
  90 came from A:
  \(P(\text{A true} \mid \text{A selected}) = 90/130 = 0.69\).
\end{itemize}

\subsubsection{5. Identifiability}

\begin{itemize}
\tightlist
\item
  \textbf{Structurally non-identifiable:} different parameter values
  give \emph{exactly} the same predictions, for any data. Examples:
  label switching in the two-trace model (Week 3), or the SDT location
  and scale (Week 2). The fix is a constraint or a reparameterization.
\item
  \textbf{Practically non-identifiable:} parameters are identifiable in
  principle but not with the available data. Recovery is poor and
  estimates trade off. The fix is a better design or a simpler model.
\end{itemize}

Parameter recovery is the practical test of identifiability.

\paragraph{For an engineer}

\begin{itemize}
\tightlist
\item
  \textbf{AIC and BIC} are the model-order selection criteria you may
  know from AR/ARMA modeling; \textbf{BIC} is essentially Rissanen's
  \textbf{MDL} criterion. The LRT for nested models is again the GLRT.
\item
  \textbf{Parameter recovery} is testing an estimator on synthetic
  signals with known ground truth, at the real SNR and record length,
  before trusting it on measured data. Engineers do this routinely; in
  psychology it is often skipped, which is why the course insists on it.
\item
  The \textbf{confusion matrix} is the channel transition matrix of a
  discrete channel: input = generating model, output = selected model.
  The \textbf{inversion matrix} is the matrix of \emph{a posteriori}
  probabilities \(P(\text{input} \mid \text{output})\), i.e.~what an MAP
  decoder needs, assuming equal input priors.
\item
  \textbf{Structural non-identifiability} is like an unobservable state
  or a rank-deficient system matrix: no amount of data helps.
  \textbf{Practical non-identifiability} is an ill-conditioned problem:
  identifiable in principle, but the noise is too large for the
  available record.
\end{itemize}

\paragraph{For the instructor}

\begin{itemize}
\tightlist
\item
  \textbf{Frame recovery as ``checking the equipment''.} No one would
  trust a thermometer that was never tested in water of known
  temperature. Recovery is that test for a model.
\item
  \textbf{Show one scatter plot of true vs.~recovered values} (lab Part
  C) before any theory: good recovery (points on the diagonal) next to
  bad recovery (a flat cloud). The picture explains itself.
\item
  \textbf{Teach confusion vs.~inversion with the small 2×2 example
  above}; students confuse them in the exam more than anything else this
  week. The medical-test analogy helps: the sensitivity of a test is not
  the probability that a positive patient is ill.
\item
  \textbf{Typical questions.}

  \begin{itemize}
  \tightlist
  \item
    \emph{``AIC or BIC?''} Depends on the goal (prediction vs.~finding
    the true model) and, practically, on which one recovers your models
    better at your design.
  \item
    \emph{``How many simulated datasets?''} 50--100 per condition is
    enough to see the pattern; more for publication.
  \item
    \emph{``Recovery failed, what now?''} Do not interpret that
    parameter. Change the design or simplify the model.
  \end{itemize}
\end{itemize}

\begin{center}\rule{0.5\linewidth}{0.5pt}\end{center}

\subsection{Worked example}

The four forgetting models fitted by MLE to the Week 2 data (training
NLLs from Week 6), with \(N = 400\) item responses:

{\def\LTcaptype{none} % do not increment counter
\begin{longtable}[]{@{}
  >{\raggedright\arraybackslash}p{(\linewidth - 8\tabcolsep) * \real{0.2000}}
  >{\raggedright\arraybackslash}p{(\linewidth - 8\tabcolsep) * \real{0.2000}}
  >{\raggedright\arraybackslash}p{(\linewidth - 8\tabcolsep) * \real{0.2000}}
  >{\raggedright\arraybackslash}p{(\linewidth - 8\tabcolsep) * \real{0.2000}}
  >{\raggedright\arraybackslash}p{(\linewidth - 8\tabcolsep) * \real{0.2000}}@{}}
\toprule\noalign{}
\begin{minipage}[b]{\linewidth}\raggedright
Model
\end{minipage} & \begin{minipage}[b]{\linewidth}\raggedright
\(k\)
\end{minipage} & \begin{minipage}[b]{\linewidth}\raggedright
NLL
\end{minipage} & \begin{minipage}[b]{\linewidth}\raggedright
AIC \(= 2\,\mathrm{NLL} + 2k\)
\end{minipage} & \begin{minipage}[b]{\linewidth}\raggedright
BIC \(= 2\,\mathrm{NLL} + k\ln 400\)
\end{minipage} \\
\midrule\noalign{}
\endhead
\bottomrule\noalign{}
\endlastfoot
exp & 2 & 29.812 & 63.62 & 71.61 \\
pow & 2 & 19.830 & \textbf{43.66} & \textbf{51.64} \\
expasym & 3 & 18.919 & 43.84 & 55.81 \\
two & 4 & 18.919 & 45.84 & 61.80 \\
\end{longtable}
}

\textbf{Step 1: AIC for the power model.}
\(2 \times 19.830 + 2 \times 2 = 43.66\).

\textbf{Step 2: BIC for the power model.} \(\ln 400 = 5.991\), so
\(2 \times 19.830 + 2 \times 5.991 = 51.64\).

\textbf{Step 3: Akaike weights.} Differences from the best AIC:
\(\Delta = (19.96, 0, 0.18, 2.18)\). Then
\(e^{-\Delta/2} = (0.00005, 1, 0.914, 0.336)\), which sum to 2.250, so

\[
w^{\text{AIC}} = (0.000,\ 0.444,\ 0.406,\ 0.149).
\]

BIC weights computed the same way: \((0.000,\ 0.884,\ 0.110,\ 0.005)\).

\textbf{Step 4: likelihood-ratio test}, exp nested in expasym:
\(\Lambda = 2(29.812 - 18.919) = 21.79\) on 1 df,
\(p \approx 3 \times 10^{-6}\). Adding an asymptote improves the
exponential model significantly.

\textbf{Interpretation.} Both criteria reject the pure exponential
decisively. AIC cannot separate the power and asymptote models (weights
0.44 vs.~0.41): for \emph{prediction} they are about equally good,
consistent with the Week 6 CV results. BIC's larger penalty favors the
two-parameter power model (0.88). The two-trace model's extra parameter
buys nothing: identical NLL to expasym, so it pays the penalty for
nothing. Whether we should \emph{trust} these preferences at this design
is a question for model recovery (lab Part D).

\begin{center}\rule{0.5\linewidth}{0.5pt}\end{center}

\subsection{Lab (2 hours)}

\subsubsection{Goals}

\begin{enumerate}
\def\labelenumi{\arabic{enumi}.}
\tightlist
\item
  Run a likelihood-ratio test and compute AIC, AICc, BIC and weights for
  the forgetting models.
\item
  Run a parameter-recovery study for the power and
  exponential-with-asymptote models, and check for trade-offs.
\item
  Run a model-recovery study for three models with AIC and BIC; build
  confusion and inversion matrices.
\item
  \textbf{Interpret}: what can and cannot be concluded from the Week 2
  data at this design?
\end{enumerate}

\textbf{Timing:} Part A 15 min · Part B 15 min · Part C 35 min · Part D
40 min · Part E 15 min.

\subsubsection{Setup}

\begin{quote}
\textbf{Code · \RL{\CMfa کد:}} \nolinkurl{code/week07.py}, cell 1 · \RL{\CMfa سلول}
1
\end{quote}

\subsubsection{Part A: likelihood-ratio test (15 min)}

\begin{quote}
\textbf{Code · \RL{\CMfa کد:}} \nolinkurl{code/week07.py}, cell 2 · \RL{\CMfa سلول}
2
\end{quote}

\textbf{Expected output:}

\Needspace{5\baselineskip}
\begin{Verbatim}
NLL exp = 29.812, NLL expasym = 18.919
LR statistic = 21.79, df = 1, naive p = 3.05e-06, boundary-corrected p = 1.52e-06
\end{Verbatim}

\textbf{Question:} why can't you use the same test to compare the power
and exponential models?

\subsubsection{Part B: information criteria and weights (15 min)}

\begin{quote}
\textbf{Code · \RL{\CMfa کد:}} \nolinkurl{code/week07.py}, cell 3 · \RL{\CMfa سلول}
3
\end{quote}

\textbf{Expected output:}

\Needspace{9\baselineskip}
\begin{Verbatim}
         k     NLL     AIC    AICc     BIC  w_AIC  w_AICc  w_BIC
model
exp      2  29.812  63.625  63.655  71.608  0.000   0.000  0.000
pow      2  19.830  43.661  43.691  51.644  0.444   0.449  0.884
expasym  3  18.919  43.837  43.898  55.812  0.406   0.405  0.110
two      4  18.919  45.837  45.939  61.803  0.150   0.146  0.006
\end{Verbatim}

\textbf{Sensitivity check:} recompute BIC with \(N = 8\) (the number of
cells) instead of 400. Does the ranking change? Which choice of \(N\) is
defensible here, and why?

\subsubsection{Part C: parameter recovery (35 min)}

\begin{quote}
\textbf{Code · \RL{\CMfa کد:}} \nolinkurl{code/week07.py}, cell 4 · \RL{\CMfa سلول}
4
\end{quote}

\textbf{Expected output} (about 1 minute):

\Needspace{15\baselineskip}
\begin{Verbatim}
pow:
   r(true, rec)   bias   RMSE
A         0.888  0.005  0.055
d         0.965  0.003  0.070

expasym:
     r(true, rec)   bias   RMSE
A           0.802  0.008  0.077
lam         0.799  0.034  0.230
s           0.918 -0.008  0.060

expasym: corr(true lam, true s) = -0.08;  corr(recovered lam, recovered s) = 0.03
\end{Verbatim}

\textbf{Interpret:} 1. Which parameters are well recovered at this
design? Which would you hesitate to interpret? 2. The expasym decay rate
\(\lambda\) has by far the largest RMSE (0.23 on a range of 0.05--1).
Look at its panel: for which true values does recovery break down, and
why? (Hint: if decay is very fast, most of the decline happens before
the second test interval, so what do the data look like?) Here the
recovered \(\lambda\) and \(s\) are \emph{not} correlated (r = 0.03), so
this is imprecision rather than a trade-off. Contrast this with the
clear \(A\)--\(d\) trade-off of the power model found in Week 5. 3. What
design change would improve recovery of \(\lambda\) and \(s\)? Test one
idea in Exercise 3.

\subsubsection{Part D: model recovery (40 min)}

\begin{quote}
\textbf{Code · \RL{\CMfa کد:}} \nolinkurl{code/week07.py}, cell 5 · \RL{\CMfa سلول}
5
\end{quote}

\textbf{Expected output} (about 1--2 minutes):

\Needspace{24\baselineskip}
\begin{Verbatim}
AIC confusion matrix (rows = generating model, columns = selected model):
          exp   pow  expasym
exp      0.91  0.00     0.09
pow      0.03  0.82     0.15
expasym  0.06  0.25     0.69
AIC inversion matrix (rows = generating model, columns = selected model):
          exp   pow  expasym
exp      0.91  0.00     0.10
pow      0.03  0.77     0.16
expasym  0.06  0.23     0.74

BIC confusion matrix (rows = generating model, columns = selected model):
          exp   pow  expasym
exp      0.99  0.01     0.00
pow      0.05  0.94     0.01
expasym  0.09  0.58     0.33
BIC inversion matrix (rows = generating model, columns = selected model):
          exp   pow  expasym
exp      0.88  0.01     0.00
pow      0.04  0.61     0.03
expasym  0.08  0.38     0.97
\end{Verbatim}

\textbf{Interpret:} 1. Which model is most often confused with which?
(With AIC, 25\% of expasym datasets are labeled power; with BIC, 58\%.)
Explain the mimicry in terms of the curves' shapes. 2. Compare AIC and
BIC. Which recovers the simpler models better, and which the more
complex one? 3. In Part B, BIC selected the power model for the real
data. Using the BIC \textbf{inversion} matrix, how confident should we
be that the data were generated by a power model? (Only 61\% of BIC's
``power'' verdicts came from power-generated data; 38\% came from
expasym.)

\subsubsection{Part E: conclusions about the Week 2 data (15 min)}

Write a paragraph (6--10 sentences) that a reviewer would accept,
summarizing what the Week 2 data say about the form of forgetting. Use
Parts A--D and your Week 6 CV results, and state what \emph{cannot} be
concluded at this design.

\subsubsection{Exercises (increasing difficulty)}

\textbf{Exercise 1 \ensuremath{\star}.} Compute AIC and BIC weights for the worked
example by hand (calculator) and check them against Part B.

\textbf{Exercise 2 \ensuremath{\star}.} Repeat Part B with AICc. For which \(N\) would
AICc and AIC differ noticeably for the two-trace model?

\textbf{Exercise 3 \ensuremath{\star}\ensuremath{\star} (design for recovery).} Rerun the expasym recovery
of Part C with intervals extended to
\(\{0.25, 1, 2, 4, 7, 14, 30, 60, 90\}\) days, and separately with
\(n = 200\) items. Which change helps \(\lambda\) and \(s\) recovery
more?

\textbf{Exercise 4 \ensuremath{\star}\ensuremath{\star} (CV in model recovery).} Replace AIC/BIC in Part D
by 5-fold CV (from Week 6) for 30 simulations per model. Is CV better or
worse than AIC at recovering the models? Is it worth the computation
time?

\textbf{Exercise 5 \ensuremath{\star}\ensuremath{\star}\ensuremath{\star} (Advanced: SDT recovery).} For the payoff
experiment of Week 4 (three conditions, 80 + 80 trials per condition),
run a model recovery for ``shared \(d'\)'' vs.~``separate \(d'\)'' using
the likelihood-ratio test at \(\alpha = .05\) and BIC. How often does
each method wrongly claim that payoffs change sensitivity? How many
trials per condition would you need to detect a true sensitivity
difference of 0.3 in 80\% of simulations?

\subsubsection{Lab check-off}

Show: (1) your IC table with weights (Part B), (2) the recovery scatter
plots with one sentence on which parameter is least well recovered and
why, (3) the AIC and BIC confusion matrices, and (4) your Part E
paragraph.

\begin{center}\rule{0.5\linewidth}{0.5pt}\end{center}

\subsection{Connection to AI \& LLMs}

\textbf{Comparing learning curves of an AI model.} LLMs can learn a new
task from a few examples in the prompt (``in-context learning''). How
the error falls with the number of examples \(k\) is a learning curve,
and the question ``exponential or power?'' is the same question as in
this week's lab (and Week 2's forgetting curves). Information criteria,
model recovery and identifiability apply unchanged. The design matters
as much as before: with few values of \(k\) and few test items, the two
curves may be impossible to tell apart.

\begin{quote}
\textbf{Code · \RL{\CMfa کد:}} \nolinkurl{code/week07.py}, cell 6 · \RL{\CMfa سلول}
6
\end{quote}

\textbf{Expected output:}

\Needspace{4\baselineskip}
\begin{Verbatim}
{'exponential': 41.6, 'power': 40.4} Akaike weights: [0.362 0.638]
\end{Verbatim}

The data were generated by a power law, but the evidence is weak
(weights 0.64 vs.~0.36).

\textbf{Try it (no API needed).} - \ensuremath{\star} Run a model recovery (Week 7, Part
D) at this design: 50 datasets from each curve. Build the confusion
matrix. - \ensuremath{\star}\ensuremath{\star} Change the design: more values of \(k\) (up to 128) or 200
test items per \(k\). Which change improves model recovery more? - \ensuremath{\star}\ensuremath{\star}
Discussion: a paper concludes from one fitted curve that ``in-context
learning follows a power law, like human practice''. What would you need
to see before accepting the comparison with humans?

\begin{center}\rule{0.5\linewidth}{0.5pt}\end{center}

\subsection{Common mistakes \& misconceptions}

\begin{enumerate}
\def\labelenumi{\arabic{enumi}.}
\tightlist
\item
  \textbf{Interpreting an AIC value on its own.} Only differences
  between models fitted to the same data mean anything.
\item
  \textbf{Using LRTs for non-nested models,} or reading a
  non-significant LRT as evidence \emph{for} the simpler model.
\item
  \textbf{Thinking BIC is ``Bayesian AIC'' with the same goal.} BIC aims
  at identifying the data-generating model; AIC aims at prediction. They
  can disagree for good reasons.
\item
  \textbf{Unclear \(N\).} Always report what counted as an observation;
  BIC and AICc depend on it.
\item
  \textbf{Skipping recovery.} A parameter that cannot be recovered at
  your design must not be interpreted, no matter how good the fit.
\item
  \textbf{Running recovery at a different design or with a different
  fitting procedure} than the real analysis. Recovery results are
  design-specific.
\item
  \textbf{Reading the confusion matrix as if it were the inversion
  matrix.} ``Data from model A are recovered as A 90\% of the time'' is
  not the same as ``when A wins, it is the true model 90\% of the
  time''.
\item
  \textbf{Choosing parameter ranges for recovery that are wider or
  narrower than plausible values.} Too wide inflates correlations; too
  narrow hides problems.
\end{enumerate}

\begin{center}\rule{0.5\linewidth}{0.5pt}\end{center}

\subsection{Check your understanding}

\textbf{1.} Write the formulas for AIC and BIC. For which \(N\) does BIC
penalize each parameter more than AIC? \textbf{Answer}
\(\mathrm{AIC} = -2\hat\ell + 2k\);
\(\mathrm{BIC} = -2\hat\ell + k\ln N\). BIC's penalty exceeds AIC's when
\(\ln N > 2\), i.e., \(N > e^2 \approx 7.4\), so for \(N \ge 8\).

\textbf{2.} Model A has AIC = 210.4, model B has AIC = 212.9. Compute
the Akaike weights and interpret them. \textbf{Answer}
\(\Delta = (0, 2.5)\), so \(e^{-\Delta/2} = (1, 0.287)\) and
\(w = (0.777, 0.223)\). A is favored about 3.5:1 as the better
predictor, but B retains non-negligible support.

\textbf{3.} Why can the likelihood-ratio test compare exponential and
exponential-with-asymptote but not exponential and power?
\textbf{Answer} The exponential is a special case of the asymptote model
(fix \(s = 0\)), so the models are nested and \(2\Delta\ell\) has a
known \(\chi^2\) distribution under the null. Exponential and power are
not nested: neither is a special case of the other.

\textbf{4.} Describe the steps of a parameter-recovery study.
\textbf{Answer} Choose plausible true parameter values; simulate data at
the real design; fit with the real fitting procedure; compare recovered
with true values (scatter plots, correlation, bias, RMSE); check
correlations among the recovered parameters for trade-offs.

\textbf{5.} In a recovery study, the recovered \(\lambda\) and \(s\)
correlate at \(r = -0.8\), although the true values were drawn
independently. What does that mean? \textbf{Answer} A trade-off: at this
design, different combinations of decay rate and asymptote produce
nearly the same predictions, so errors in one are compensated by the
other. Individual values of \(\lambda\) and \(s\) should not be
interpreted separately. A better design or a reparameterization is
needed.

\textbf{6.} Explain the difference between a confusion matrix and an
inversion matrix in model recovery. \textbf{Answer} Confusion:
\(P(\text{selected} \mid \text{generated})\), with rows summing to 1,
which describes the method's behavior for each true model. Inversion:
\(P(\text{generated} \mid \text{selected})\), with columns summing to 1
(assuming equal priors for the generating models), which tells us how
much to trust a given winner when analyzing real data.

\textbf{7.} A model recovery shows that AIC selects the three-parameter
model 40\% of the time when the two-parameter model generated the data.
What would you do? \textbf{Answer} Treat AIC wins for the complex model
with caution at this design: check BIC or CV recovery, which may be more
conservative, improve the design (more trials, more informative
conditions), and report the recovery results alongside the real-data
comparison.

\begin{center}\rule{0.5\linewidth}{0.5pt}\end{center}

\subsection{Assignment}

\subsubsection{A3 continues (due at the start of the Week 9 lecture)}

Tasks 4 (information criteria) and 5 (parameter and model recovery) of
A3 use this week's material. The full brief is in Week 6. Reuse today's
\nolinkurl{parameter_recovery} code and the model-recovery loop: they
are designed to be adapted directly.

\subsubsection{Final project: teams and proposal brief released}

\begin{itemize}
\tightlist
\item
  \textbf{Teams:} 1--2 students. Register your team with the TA by the
  end of this week.
\item
  \textbf{Task:} choose a cognitive model and apply the full pipeline to
  a dataset (simulated, open, or your own): fit, evaluate (at least one
  model comparison with an alternative model), recover (at least one
  parameter-recovery analysis), interpret. \textbf{Pairs} also include a
  Bayesian/MCMC or hierarchical analysis and an individual contribution
  statement.
\item
  \textbf{Proposal (due at the start of the Week 9 lecture, max. 2
  pages):} research question, model(s) with equations, data source and
  design, planned comparison and recovery analyses, and a timeline. The
  full project brief and rubric are distributed separately
  (\nolinkurl{project_brief.md}).
\item
  \textbf{Not sure what to model?} The topics of Weeks 14--15 (exemplar
  categorization, drift-diffusion, reinforcement learning) all make good
  projects, as do extensions of the forgetting and SDT models from this
  course.
\end{itemize}

\begin{center}\rule{0.5\linewidth}{0.5pt}\end{center}

\subsection{Further reading}

\begin{itemize}
\tightlist
\item
  \textbf{Farrell \& Lewandowsky (2018), chapter 10.} Model comparison:
  likelihood-ratio tests, AIC, BIC and their weights; their assumptions
  and limitations. \emph{(Main reading.)}
\item
  \textbf{Lee \& Wagenmakers (2014), chapter 7} (introduction). How the
  Bayesian approach to model comparison relates to the methods of this
  week, as preparation for Week 10.
\item
  \emph{Optional, outside the course list:} Wilson, R. C., \& Collins,
  A. G. E. (2019). Ten simple rules for the computational modeling of
  behavioral data. \emph{eLife, 8}, e49547. A practical guide that
  treats parameter and model recovery as standard practice.
\end{itemize}

\emph{Preview for Week 8:} the midterm. The lecture is an integrated
review in which we run the whole pipeline (simulate, fit, evaluate,
recover, interpret) on a new dataset; the lab slot is the midterm exam.
